% Respuesta a la pregunta 5
Las bases de datos NoSQL son bases de datos no relacionales, es decir, no almacenan datos en tablas relacionales basadas en reglas. Estás son más flexibles al admitir datos no estructurados como documentos, pares clave-valor y gráficos. Hay distintas bases de datos NoSQL, son útiles según el contexto. Las principales son: Bases de datos de documentos, Bases de datos de pares clave-valor, Bases de datos orientadas a columnas, Bases de datos de grafos y Base de Datos en Memoria (Google Cloud, s.f). 

Su principal ventaja está en su principal característica, la flexibilidad de admitir datos semi estructurados o no estructurados. Al permitir esa flexibilidad en la estructura de datos, permite que el desarrollo de una app sea más ágil, pues los desarrolladores dedican menos tiempo a la transformación de datos, agilizando el desarrollo (Google Cloud, s.f).
IBM (s.f) señala que son más rentables, pues las bases de datos NoSQL le permiten escalar de manera rápida y horizontal, asignando mejor los recursos para minimizar los costos.

A comparación de SQL, NoSQL es reciente, por está razón los desarrolladores tienen menos experiencia, menos herramientas y productos disponibles y menos ayuda si surgen problemas sin documentar (Google Cloud, s.f). 
En varios casos, las bases de datos NoSQL a comparación con las bases de datos relacionales, no cuentan con la seguridad de los datos y un alto nivel de coherencia de datos que son estándar en las bases de datos SQL (Google Cloud, s.f). 
Generalmente, no son una buena opción para las aplicaciones que ejecutan consultas y uniones complejas (Google Cloud, s.f).
